\chapter{高等数学：函数极限}

\section{函数极限}

\begin{problemblock}
\textbf{例1.} 已知 \(\displaystyle \lim_{x\to+\infty}f(x)\) 存在，且
\[
f(x)=\frac{x^{1+x}}{(1+x)^x}-\frac{x}{e}
  +2\lim_{x\to+\infty}f(x),
\]
求 \(f(x)\)。
\end{problemblock}

\begin{solutionblock}
\analysis{本题是“含未知极限的函数表达式”。考研中遇到这种题，先设未知极限，再对等式两端取极限。难点在于精确求
\(\frac{x^{1+x}}{(1+x)^x}-\frac{x}{e}\) 的极限。}

设
\[
L=\lim_{x\to+\infty}f(x).
\]
先处理主项：
\[
\frac{x^{1+x}}{(1+x)^x}
=x\left(\frac{x}{x+1}\right)^x
=x\left(1+\frac1x\right)^{-x}.
\]
因此需要求
\[
\lim_{x\to+\infty}x\left[\left(1+\frac1x\right)^{-x}-\frac1e\right].
\]
令
\[
u_x=\left(1+\frac1x\right)^{-x}.
\]
由于
\[
\ln u_x=-x\ln\left(1+\frac1x\right)
=-x\left(\frac1x-\frac1{2x^2}+O\left(\frac1{x^3}\right)\right)
=-1+\frac1{2x}+O\left(\frac1{x^2}\right),
\]
所以
\[
u_x=e^{-1}\exp\left(\frac1{2x}+O\left(\frac1{x^2}\right)\right)
=\frac1e\left(1+\frac1{2x}+O\left(\frac1{x^2}\right)\right).
\]
于是
\[
\frac{x^{1+x}}{(1+x)^x}
=\frac{x}{e}+\frac1{2e}+O\left(\frac1x\right),
\]
从而
\[
\lim_{x\to+\infty}\left(\frac{x^{1+x}}{(1+x)^x}-\frac{x}{e}\right)=\frac1{2e}.
\]
对原式取极限得
\[
L=\frac1{2e}+2L,
\]
故
\[
L=-\frac1{2e}.
\]
代回原式，
\[
f(x)=\frac{x^{1+x}}{(1+x)^x}-\frac{x}{e}-\frac1e.
\]

\method{方法二：洛必达型处理}
也可把
\[
x\left[\left(1+\frac1x\right)^{-x}-\frac1e\right]
\]
看作 \(\infty\cdot0\) 型，令 \(t=1/x\)，转为
\[
\lim_{t\to0^+}\frac{(1+t)^{-1/t}-e^{-1}}{t},
\]
这就是函数 \(g(t)=(1+t)^{-1/t}\) 在 \(0\) 处的右导数。由展开或求导可得该极限为 \(1/(2e)\)。

\answer{\(\displaystyle f(x)=\frac{x^{1+x}}{(1+x)^x}-\frac{x}{e}-\frac1e\)。}
\examnote{这题关键不是套等价无穷小，而是保留到常数项；只知道 \((1+1/x)^x\to e\) 不够。}
\end{solutionblock}

\begin{problemblock}
\textbf{例2.} 已知 \(\varphi(x)\) 是 \(x\to0\) 时的无穷小量，则下列说法中正确的有几项？
\[
\begin{array}{ll}
\text{① } \displaystyle\lim_{x\to0}\frac{\sin\varphi(x)}{\varphi(x)}=1; &
\text{② 当 }x\to0\text{ 时，}\sin\varphi(x)\text{ 与 }\varphi(x)\text{ 是等价无穷小；}\\[2mm]
\text{③ 当 }x\to0\text{ 时，}\displaystyle\frac1{\varphi(x)}\text{ 是无穷大量；}&
\text{④ 若连续函数 }f(x)\text{ 满足 }\displaystyle\lim_{x\to0}f(x)=1,\\
&\qquad\text{则 }\displaystyle\lim_{x\to0}f(\varphi(x))=1.
\end{array}
\]
A. 0 \qquad B. 1 \qquad C. 2 \qquad D. 3
\end{problemblock}

\begin{solutionblock}
\analysis{本题考查复合极限与无穷小的基本性质。容易误判的是③：无穷小的倒数要成为无穷大量，通常还需保证该无穷小在去心邻域内不为零。题干只说 \(\varphi(x)\to0\)，没有说明 \(\varphi(x)\ne0\)。}

① 因为 \(\varphi(x)\to0\)，令 \(u=\varphi(x)\)，由基本极限
\[
\lim_{u\to0}\frac{\sin u}{u}=1
\]
可得①正确。

② 由①直接得到
\[
\sin\varphi(x)\sim\varphi(x),
\]
故②正确。

③ 不一定正确。例如取
\[
\varphi(x)\equiv0,
\]
它当然是 \(x\to0\) 时的无穷小，但 \(\frac1{\varphi(x)}\) 根本无定义，不能称为无穷大量。即使取在某些点为零的无穷小，也会破坏该说法。因此③错误。

④ 因为 \(\varphi(x)\to0\)，且 \(\lim_{u\to0}f(u)=1\)，复合极限定理给出
\[
\lim_{x\to0}f(\varphi(x))=1.
\]
故④正确。

正确的是①②④，共 3 项。

\answer{D}
\examnote{考研选择题里“无穷小的倒数是无穷大”一定要看是否默认非零；题干没给非零条件时要谨慎。}
\end{solutionblock}

\begin{problemblock}
\textbf{例3.} 当 \(x\to0\) 时，\(\alpha(x),\beta(x)\) 是非零无穷小量，则以下命题中：
\[
\begin{array}{ll}
\text{① }\alpha(x)\sim\beta(x),\text{ 则 }\alpha^2(x)\sim\beta^2(x);&
\text{② }\alpha^2(x)\sim\beta^2(x),\text{ 则 }\alpha(x)\sim\beta(x);\\
\text{③ }\alpha(x)\sim\beta(x),\text{ 则 }\alpha(x)-\beta(x)=o(\alpha(x));&
\text{④ }\alpha(x)-\beta(x)=o(\alpha(x)),\text{ 则 }\alpha(x)\sim\beta(x).
\end{array}
\]
所有真命题的序号是（\quad）

A. ①③ \qquad B. ①④ \qquad C. ①③④ \qquad D. ②③④
\end{problemblock}

\begin{solutionblock}
\analysis{本题考查等价无穷小的定义。核心是把所有命题都化成比值极限。}

① 若 \(\alpha\sim\beta\)，则
\[
\frac{\alpha^2}{\beta^2}=\left(\frac{\alpha}{\beta}\right)^2\to1,
\]
所以①正确。

② 错误。取
\[
\alpha(x)=x,\qquad \beta(x)=-x.
\]
则
\[
\alpha^2(x)=\beta^2(x),
\]
因此 \(\alpha^2\sim\beta^2\)，但
\[
\frac{\alpha(x)}{\beta(x)}=-1\not\to1,
\]
所以 \(\alpha\not\sim\beta\)。

③ 若 \(\alpha\sim\beta\)，则
\[
\frac{\alpha-\beta}{\alpha}=1-\frac{\beta}{\alpha}\to0,
\]
所以
\[
\alpha-\beta=o(\alpha).
\]
③正确。

④ 若
\[
\alpha-\beta=o(\alpha),
\]
则
\[
\frac{\alpha-\beta}{\alpha}\to0,
\]
即
\[
1-\frac{\beta}{\alpha}\to0.
\]
于是
\[
\frac{\beta}{\alpha}\to1,
\]
从而 \(\alpha\sim\beta\)。④正确。

\answer{C}
\examnote{“平方等价”只能推出绝对值量级相同，不能推出符号一致，这是②的陷阱。}
\end{solutionblock}

\begin{problemblock}
\textbf{例4.} 当 \(x\to0\) 时，下列无穷小量最高阶的是（\quad）
\[
\begin{array}{ll}
\text{A. }\displaystyle\int_0^{5x}(1+t)^{1/t}\,dt,&
\text{B. }\displaystyle\int_{x^2}^{1-\cos x}\frac{\sin t}{t}\,dt,\\[2mm]
\text{C. }\displaystyle\int_x^{1-\cos x}\ln(1+t^2)\,dt,&
\text{D. }\displaystyle\int_x^{\sin x}(e^{x^2}-1)\,dt.
\end{array}
\]
\end{problemblock}

\begin{solutionblock}
\analysis{比较无穷小阶数，核心是求主部。含变限积分时，先看被积函数主部，再看上下限差的阶数。}

A 中
\[
(1+t)^{1/t}\to e,
\]
所以
\[
\int_0^{5x}(1+t)^{1/t}\,dt\sim 5ex,
\]
是一阶无穷小。

B 中
\[
\frac{\sin t}{t}\to1,\qquad 1-\cos x\sim\frac{x^2}{2},
\]
故
\[
\int_{x^2}^{1-\cos x}\frac{\sin t}{t}\,dt
\sim (1-\cos x)-x^2
\sim -\frac{x^2}{2},
\]
是二阶无穷小。

C 中
\[
\ln(1+t^2)\sim t^2.
\]
又 \(1-\cos x=O(x^2)\)，所以
\[
\int_x^{1-\cos x}\ln(1+t^2)\,dt
\sim\int_x^0t^2\,dt=-\frac{x^3}{3},
\]
是三阶无穷小。

D 中注意被积函数 \(e^{x^2}-1\) 与积分变量 \(t\) 无关，因此
\[
\int_x^{\sin x}(e^{x^2}-1)\,dt
=(e^{x^2}-1)(\sin x-x).
\]
而
\[
e^{x^2}-1\sim x^2,\qquad \sin x-x\sim-\frac{x^3}{6},
\]
所以
\[
\int_x^{\sin x}(e^{x^2}-1)\,dt
\sim-\frac{x^5}{6},
\]
是五阶无穷小。

最高阶为 D。

\answer{D}
\examnote{D 的被积函数写的是 \(e^{x^2}-1\)，不是 \(e^{t^2}-1\)。考场上这类细节决定答案。}
\end{solutionblock}

\begin{problemblock}
\textbf{例5.} 已知
\[
\lim_{x\to0}\frac{2-2\cos x+x f(x)}{x^3}=0,
\]
计算
\[
(1)\ \lim_{x\to0}\frac{x+f(x)}{x^2};\qquad
(2)\ \lim_{x\to0}\frac{x+f(\sin x)}{x^2}.
\]
\end{problemblock}

\begin{solutionblock}
\analysis{本题由极限条件反推 \(f(x)\) 的局部主部。要注意第二问是 \(f(\sin x)\)，不能直接把第一问结果机械代入。}

由题意
\[
2-2\cos x+x f(x)=o(x^3).
\]
又
\[
2-2\cos x=x^2-\frac{x^4}{12}+O(x^6),
\]
所以
\[
x f(x)=-(2-2\cos x)+o(x^3),
\]
即
\[
f(x)=-\frac{2-2\cos x}{x}+o(x^2).
\]
由于
\[
\frac{2-2\cos x}{x}=x-\frac{x^3}{12}+O(x^5),
\]
因此
\[
f(x)=-x+o(x^2).
\]
于是
\[
\frac{x+f(x)}{x^2}\to0.
\]
故第一问答案为 0。

第二问中令 \(u=\sin x\)，由上式对任意 \(u\to0\) 成立，
\[
f(u)=-\frac{2-2\cos u}{u}+o(u^2).
\]
于是
\[
f(\sin x)=-\frac{2-2\cos(\sin x)}{\sin x}+o(\sin^2x).
\]
利用
\[
\frac{2-2\cos u}{u}=u-\frac{u^3}{12}+O(u^5),
\]
可得
\[
f(\sin x)=-\sin x+O(x^3)+o(x^2).
\]
所以
\[
x+f(\sin x)=x-\sin x+O(x^3)+o(x^2)=o(x^2),
\]
从而
\[
\lim_{x\to0}\frac{x+f(\sin x)}{x^2}=0.
\]

\method{方法二：只抓阶数}
由已知直接得到 \(f(x)+x=o(x^2)\)。对第二问，用 \(u=\sin x\) 得 \(f(\sin x)+\sin x=o(x^2)\)，于是
\[
x+f(\sin x)=x-\sin x+o(x^2)=o(x^2).
\]
这比完整展开更快。

\answer{(1) \(0\)；(2) \(0\)。}
\examnote{若已推出 \(f(x)+x=o(x^2)\)，第二问要补上 \(x-\sin x=O(x^3)\)，这样才能严谨说明除以 \(x^2\) 后趋零。}
\end{solutionblock}

\begin{problemblock}
\textbf{例6.} 若函数
\[
f(x)=\frac{e^x-b}{(x-a)(x-b)}
\]
有可去间断点，则 \(ab\) 的取值范围为（\quad）

A. \([e,+\infty)\) \qquad
B. \(\left[-\frac1e,e\right]\) \qquad
C. \(\left(-\frac1e,e\right)\) \qquad
D. \(\left[-\frac1e,+\infty\right)\)
\end{problemblock}

\begin{solutionblock}
\analysis{可去间断点来自“分母为零、分子也为零，并能约去”。本题要把条件转化为 \(a,b\) 的关系，再求 \(ab\) 的范围。}

分母零点为 \(x=a\) 与 \(x=b\)。若在某点 \(c\) 可去间断，则必须有
\[
e^c-b=0.
\]
也就是
\[
c=\ln b,\qquad b>0.
\]

若可去点为 \(x=a\)，则
\[
a=\ln b.
\]
此时
\[
ab=b\ln b,\qquad b>0.
\]
令
\[
\phi(b)=b\ln b.
\]
则
\[
\phi'(b)=\ln b+1.
\]
当 \(b=e^{-1}\) 时取最小值
\[
\phi(e^{-1})=-\frac1e.
\]
且当 \(b\to+\infty\) 时，\(b\ln b\to+\infty\)。因此
\[
ab\in\left[-\frac1e,+\infty\right).
\]

若可去点为 \(x=b\)，则需
\[
e^b-b=0,
\]
但 \(e^b>b\) 对一切实数 \(b\) 成立，因此无解。

\answer{D}
\examnote{不要把“分母有两个零点”误认为两个都能约去；\(e^b=b\) 在实数范围无解。}
\end{solutionblock}

\begin{problemblock}
\textbf{例7.} 设函数 \(y=f(x)\) 由参数方程
\[
\begin{cases}
x=\dfrac1{t-2},\\[1mm]
y=\dfrac{t\ln|t|}{|t-t^3|}
\end{cases}
\]
确定，则 \(f(x)\) 的第一类间断点的个数为（\quad）

A. 0 \qquad B. 1 \qquad C. 2 \qquad D. 3
\end{problemblock}

\begin{solutionblock}
\analysis{参数方程确定函数时，先找参数表达式中可能出问题的 \(t\)，再通过 \(x=1/(t-2)\) 映射到 \(x\)-点。第一类间断点要求左右极限都存在且有限。}

当 \(t\ne0,\pm1,2\) 时，
\[
y=\frac{t\ln|t|}{|t-t^3|}
=\frac{\ln|t|}{|1-t^2|}.
\]
可能产生间断的参数为
\[
t=0,\quad t=1,\quad t=-1.
\]
其中 \(t=2\) 对应 \(x\to\infty\)，不对应有限的 \(x\) 间断点。

当 \(t\to1\) 时，
\[
\ln t\sim t-1,\qquad |1-t^2|\sim2|t-1|.
\]
所以
\[
\lim_{t\to1^+}y=\frac12,\qquad
\lim_{t\to1^-}y=-\frac12.
\]
两侧极限有限但不相等，对应
\[
x=\frac1{1-2}=-1,
\]
是第一类间断点。

当 \(t\to-1\) 时，令 \(t=-1+s\)，则
\[
\ln|t|\sim -s,\qquad |1-t^2|\sim2|s|.
\]
故两侧极限分别为 \(\pm\frac12\)，有限但不相等，对应
\[
x=\frac1{-1-2}=-\frac13,
\]
也是第一类间断点。

当 \(t\to0\) 时，
\[
y=\frac{\ln|t|}{|1-t^2|}\to-\infty,
\]
对应
\[
x=\frac1{0-2}=-\frac12.
\]
该点为无穷间断点，不是第一类间断点。

所以第一类间断点共有 2 个。

\answer{C}
\examnote{第一类间断点必须左右极限有限；趋于无穷的不算第一类。}
\end{solutionblock}

\begin{problemblock}
\textbf{例8.} 设 \(f(x)\) 在 \((0,+\infty)\) 内有定义，有一个间断点，而
\[
g(x)=\lim_{n\to\infty}\frac{\ln(e^n+x^n)}{n}\qquad(x>0),
\]
则（\quad）

A. \(g[f(x)]\) 必有间断点。\qquad
B. \(f[g(x)]\) 必有间断点。

C. \([f(x)]^2\) 必有间断点。\qquad
D. \(\dfrac{f(x)}{g(x)}\) 必有间断点。
\end{problemblock}

\begin{solutionblock}
\analysis{本题是复合函数与四则运算的间断性判断。先求出 \(g(x)\)，再判断哪些操作一定保持间断。}

用最大项原则：
\[
\lim_{n\to\infty}\frac{\ln(e^n+x^n)}{n}
=\max\{1,\ln x\}.
\]
因此
\[
g(x)=
\begin{cases}
1,&0<x\le e,\\
\ln x,&x>e.
\end{cases}
\]
所以 \(g(x)\) 在 \((0,+\infty)\) 上连续，且 \(g(x)>0\)。

对 D，若 \(\dfrac{f(x)}{g(x)}\) 在 \(f\) 的那个间断点 \(x_0\) 处连续，由于 \(g(x)\) 连续且 \(g(x_0)>0\)，则
\[
f(x)=g(x)\cdot\frac{f(x)}{g(x)}
\]
也应在 \(x_0\) 连续，这与 \(f\) 在 \(x_0\) 间断矛盾。因此 \(\dfrac{f(x)}{g(x)}\) 必有间断点。

其余选项不一定。比如平方可能消除跳跃间断：若某点附近 \(f\) 左右分别为 \(1\) 与 \(-1\)，则 \(f^2\) 可连续。复合 \(g[f(x)]\) 或 \(f[g(x)]\) 也可能因为值域落在连续区域或被 \(g\) 的常值段消除间断，不能保证。

\answer{D}
\examnote{非零连续函数作乘除不会消除原函数的间断，这是本题最稳的逻辑入口。}
\end{solutionblock}

\begin{problemblock}
\textbf{例9.} 设函数 \(f(x)\) 在 \((0,\pi)\) 内可导，\(f(x)>0\)，
\[
f\left(\frac\pi2\right)=\frac4{\pi^2},
\]
且
\[
\lim_{h\to0}\left[\frac{f(x+h\sin x)}{f(x)}\right]^{1/h}
=e^{\frac{2(x\cos x-\sin x)}{x}},\qquad x\in(0,\pi).
\]
(1) 求 \(f(x)\)；(2) 证明 \(f(x)\) 在 \((0,\pi)\) 上有界。
\end{problemblock}

\begin{solutionblock}
\analysis{这是“指数型极限反推出微分方程”。因为 \(f(x)>0\)，可以取对数。}

对极限取对数：
\[
\lim_{h\to0}\frac1h
\ln\frac{f(x+h\sin x)}{f(x)}
=\frac{2(x\cos x-\sin x)}{x}.
\]
左边是复合增量的导数。因为
\[
x+h\sin x-x=h\sin x,
\]
所以
\[
\ln f(x+h\sin x)-\ln f(x)
\sim h\sin x\cdot \frac{f'(x)}{f(x)}.
\]
故
\[
\sin x\frac{f'(x)}{f(x)}
=\frac{2(x\cos x-\sin x)}{x}.
\]
在 \(x\in(0,\pi)\) 内 \(\sin x>0\)，于是
\[
\frac{f'(x)}{f(x)}
=2\cot x-\frac2x.
\]
积分得
\[
\ln f(x)=2\ln(\sin x)-2\ln x+C,
\]
即
\[
f(x)=C\left(\frac{\sin x}{x}\right)^2.
\]
由
\[
f\left(\frac\pi2\right)=C\left(\frac{1}{\pi/2}\right)^2
=C\frac4{\pi^2}
=\frac4{\pi^2},
\]
得 \(C=1\)。因此
\[
f(x)=\left(\frac{\sin x}{x}\right)^2.
\]

下面证明有界。对 \(x\in(0,\pi)\)，有
\[
0<\sin x\le x,
\]
所以
\[
0<f(x)=\left(\frac{\sin x}{x}\right)^2\le1.
\]
故 \(f(x)\) 在 \((0,\pi)\) 上有界。

\method{一题多解提示}
也可令 \(F(x)=\ln f(x)\)，原极限直接化为
\[
\lim_{h\to0}\frac{F(x+h\sin x)-F(x)}{h}
=F'(x)\sin x,
\]
从而更快得到同一个微分方程。

\answer{\(\displaystyle f(x)=\left(\frac{\sin x}{x}\right)^2\)，且 \(0<f(x)\le1\)。}
\examnote{处理 \([\,\cdot\,]^{1/h}\) 型极限，第一反应是取对数；这是考研中最稳定的标准动作。}
\end{solutionblock}

\begin{problemblock}
\textbf{例10.} (1) 证明：当 \(x\ge1\) 时，
\[
\frac1{1+x}<\ln\left(1+\frac1x\right)<\frac1x.
\]
(2) 设函数 \(f(x)\) 在区间 \([1,+\infty)\) 上连续可导，且
\[
f'(x)=\frac1{1+f^2(x)}
\left[\sqrt{\frac1x}-\sqrt{\ln\left(1+\frac1x\right)}\right],
\]
证明：\(\displaystyle\lim_{x\to+\infty}f(x)\) 存在。
\end{problemblock}

\begin{solutionblock}
\analysis{第一问是对数不等式；第二问要把微分方程转为 \(\arctan f(x)\) 的收敛性。}

\method{(1) 积分证明}
因为
\[
\ln\left(1+\frac1x\right)=\int_x^{x+1}\frac1t\,dt.
\]
当 \(t\in[x,x+1]\) 时，
\[
\frac1{x+1}\le\frac1t\le\frac1x.
\]
且不恒等，于是积分得
\[
\frac1{x+1}<\int_x^{x+1}\frac1t\,dt<\frac1x,
\]
即
\[
\frac1{1+x}<\ln\left(1+\frac1x\right)<\frac1x.
\]

\method{(1) 函数法}
也可令
\[
\phi(u)=\ln(1+u)-\frac{u}{1+u},\qquad u>0.
\]
则
\[
\phi'(u)=\frac1{1+u}-\frac1{(1+u)^2}>0,\qquad \phi(0)=0,
\]
所以 \(\ln(1+u)>\frac{u}{1+u}\)。取 \(u=1/x\) 得左边不等式。右边由
\[
\ln(1+u)<u
\]
得到。

\method{(2) 证明极限存在}
由第一问，
\[
0<\ln\left(1+\frac1x\right)<\frac1x,
\]
所以
\[
\sqrt{\frac1x}-\sqrt{\ln\left(1+\frac1x\right)}>0.
\]
注意这里题目中已经有
\[
f'(x)=\frac1{1+f^2(x)}\left[\sqrt{\frac1x}-\sqrt{\ln\left(1+\frac1x\right)}\right].
\]
为了利用反正切，应改看 \(F(x)=\int_0^{f(x)}(1+u^2)\,du=f(x)+\frac{f^3(x)}3\)。则
\[
F'(x)=(1+f^2(x))f'(x)
=\sqrt{\frac1x}-\sqrt{\ln\left(1+\frac1x\right)}.
\]
下面证明右端在 \([1,+\infty)\) 上可积。

记
\[
A(x)=\sqrt{\frac1x}-\sqrt{\ln\left(1+\frac1x\right)}.
\]
因为
\[
A(x)=
\frac{\frac1x-\ln(1+\frac1x)}
{\sqrt{\frac1x}+\sqrt{\ln(1+\frac1x)}}.
\]
又
\[
\frac1x-\ln\left(1+\frac1x\right)
=\int_0^{1/x}\frac{t}{1+t}\,dt
<\int_0^{1/x}t\,dt
=\frac1{2x^2},
\]
并且由第一问
\[
\ln\left(1+\frac1x\right)>\frac1{x+1},
\]
故
\[
\sqrt{\frac1x}+\sqrt{\ln\left(1+\frac1x\right)}
\ge \sqrt{\ln\left(1+\frac1x\right)}
>\frac1{\sqrt{x+1}}.
\]
因此当 \(x\ge1\) 时，
\[
0<A(x)<\frac{\sqrt{x+1}}{2x^2}\le \frac{C}{x^{3/2}},
\]
其中 \(C\) 为常数。由于
\[
\int_1^{+\infty}\frac{dx}{x^{3/2}}
\]
收敛，故
\[
\int_1^{+\infty}A(x)\,dx
\]
收敛。

于是
\[
F(x)=F(1)+\int_1^x A(t)\,dt
\]
随 \(x\to+\infty\) 有极限。函数
\[
F(u)=u+\frac{u^3}{3}
\]
严格单调递增，且 \(F(u)\to\pm\infty\) 当 \(u\to\pm\infty\)。因此由 \(F(f(x))\) 有极限，可推出 \(f(x)\) 有极限。

\answer{结论成立，\(\displaystyle\lim_{x\to+\infty}f(x)\) 存在。}
\examnote{这里不能误用 \(\arctan f\)。因为题目给的是 \(f'=\frac{A(x)}{1+f^2}\)，真正一乘就消掉分母的是 \(f+\frac{f^3}{3}\)。}
\end{solutionblock}
